\section{Analysis}
\label{sec:analysis}

\paragraph{Intervention behavior.}
We further analyze the behavior of the intervention gate in \figref{fig:intervention_analysis} in \appref{sec:appendix_intervention_analysis}.
%
According to this figure, \ours{} assigns more \textsc{Skip} labels and fewer \textsc{Weak}/\textsc{Strong} labels as training progresses, indicating that the student gradually produces more acceptable actions and requires less teacher correction.
%
Across turn positions, \textsc{Skip} remains the dominant label but becomes less frequent in later turns, while deterministic failures increase, suggesting that response quality becomes harder to maintain as interaction histories grow longer.
%
These trends support the motivation for selective multi-turn OPD: \ours{} does not uniformly downweight late turns, but distinguishes acceptable behavior, uncertain behavior, and clear execution failures. Detailed label distributions and analysis are provided in \appref{sec:appendix_intervention_analysis}.

\paragraph{Ablation study.}
We conduct incremental ablations on ALFWorld to isolate the contribution of each component in \ours{}.
%
The \textit{Intervention-only} variant already performs strongly, showing that teacher-judged turn-level intervention can identify useful correction points without requiring a task-level verifier.
%
Adding teacher confidence gives \ours{}-NoNorm, which improves some settings but can also reduce unseen generalization when selective weights shrink the effective training signal.
%
The full \ours{} objective addresses this issue with vanilla-scale normalization, improving unseen SR from $64.93$ to $70.15$ for Qwen3-1.7B distilled from Qwen3-8B and from $69.40$ to $70.90$ for Qwen3-0.6B distilled from Qwen3-32B.
%
Overall, intervention weighting, teacher confidence, and loss-scale normalization play complementary roles in selective multi-turn OPD.
%
Full ablation results are provided in \appref{sec:appendix_ablation}.

\paragraph{Sensitivity Analysis}
\label{sec:sensitivity_analysis}

\begin{table}[t]
\centering
\caption{Sensitivity analysis on ALFWorld using Qwen3-8B as teacher and Qwen3-1.7B as student. The default configuration uses $\alpha=0.5$, teacher temperature $0.0$, and vanilla-scale normalization over all student tokens.}
\label{tab:sensitivity}
\resizebox{\columnwidth}{!}{%
\begin{tabular}{llcccc}
\toprule
\textbf{Type} & \textbf{Setting} & \textbf{Seen SR} & \textbf{Unseen SR} & \textbf{Seen Turns} & \textbf{Unseen Turns} \\
\midrule
\multirow{4}{*}{$\alpha$}
& $0.3$ & 63.57 & 67.16 & 17.34 & 16.78 \\
& $0.5$ \;(\textit{default}) & \textbf{64.29} & \textbf{70.15} & 18.14 & 17.86 \\
& $0.7$ & 57.86 & 66.42 & 18.86 & 18.47 \\
& $1.0$ & 58.57 & 61.94 & 19.31 & 19.02 \\
\midrule
\multirow{5}{*}{Teacher temp.}
& $0.0$ \;(\textit{default}) & \textbf{64.29} & \textbf{70.15} & 18.14 & 17.86 \\
& $0.3$ & 60.00 & 62.69 & 18.61 & 18.35 \\
& $0.7$ & 60.71 & 65.67 & 18.77 & 18.34 \\
& $1.0$ & 60.71 & 63.43 & 18.24 & 18.11 \\
& $1.5$ & 57.86 & 63.43 & 19.24 & 18.74 \\
\midrule
Normalization
& Active-token norm. & 62.14 & 68.66 & 18.41 & \textbf{17.46} \\
\bottomrule
\end{tabular}%
}
\end{table}

We analyze the sensitivity of \ours{} on ALFWorld using Qwen3-8B as teacher and Qwen3-1.7B as student.
%
As shown in \tabref{tab:sensitivity}, the default weak-intervention weight $\alpha=0.5$ achieves the best seen and unseen SR.
%
A smaller value, $\alpha=0.3$, remains competitive but underperforms on the unseen split, suggesting that weak interventions still provide useful supervision.
%
Larger values, such as $\alpha=0.7$ and $\alpha=1.0$, degrade performance and increase the average number of turns, indicating that treating weak interventions too similarly to strong interventions can over-regularize the student.
%
We also vary the teacher temperature used for confidence estimation.
%
The default deterministic setting, temperature $0.0$, performs best, while higher temperatures consistently reduce success rate, likely because the confidence score becomes less stable.
%
Finally, an alternative normalization that counts only active, non-skipped tokens remains competitive with $68.66\%$ unseen SR, but underperforms the default vanilla-scale normalization at $70.15\%$.
%
Overall, these results support using a moderate weak-intervention weight, deterministic teacher confidence, and normalization over all student tokens as the default configuration.

